\scanpage{29}

\deleted{\unsure{Il faut bien} $x,y\notin$}
\margin{\unsure{cf le NB en final}}
\figure{\begin{tikzpicture}\draw (0,0) .. controls (-.5,.8) and (1.9,1) .. (2.4,.1)
.. controls (2.6,-.25) and (1.8,-.2) .. (1.6,-.35)
.. controls (1.5,-.5) and (1.5,-.7) .. (1.5,-.8);
\draw (0,0)..controls(.6,-.2)and(1.2,-.1)..(1.8,-.15);\end{tikzpicture}}
Par construction, on a $x\in K$, $x\notin K^p$, et
$y\notin K^p$, $y\in K^p(x)$, avec \deleted{$\unsure{\alpha,\beta}$}
\inserted{les $a_i$} sont les \uline{seuls} $\in K^p$
tels qu'on ait \deleted{$y=\alpha x+\beta$. \unsure{Evident}}
\inserted{$y=\sum_0^{p-1}a_i x^i$}.

Comme $x\notin K^p$, on a $x\notin K^p\cap k$, mais bien
sur $x^p\in K^p\cap k$ ; si donc $(x,y)$ \unsure{n'étaient}
\deleted{lin.} $p$-\unsure{libres} sur $K^p\cap k$, on aurait une
relation $y=\sum_0^{p-1}a_i' x^i$, les $a_i'\in K^p\cap k$, donc
\deleted{\unsure{les}} \inserted{on aurait} $a_i'=a_i$, donc
$a_i\in k=\mathbb F(x,y)$, absurde.

N.B. Condition essentielle
\[
\left\{\begin{array}{l}
(x,y)\quad p\text{-libres sur }\lambda=\mathbb F(a_1,\ldots,a_{p-1})\\
\deleted{\unsure{a_1}}\quad\lambda\unsure{\not\subset|\notin}\mathbb F(x,y)
\end{array}\right.
\]
\deleted{Il suffit \unsure{par}}
exemple
\deleted{que $x,y$ \unsure{soient} \inserted{\unsure{$p$-indépendants}}
\unsure{sur} $\lambda=\mathbb F(a_0,\ldots,a_{p-1})$, et
\unsure{que prendre} $\lambda$ \unsure{en un autre}]
\inserted{\unsure{de} $\mathbb F(x,y,a)$}}

Prendre $x,y,a_1$ \deleted{\unsure{alg.}} $p$-indép. sur $\unsure{\mathbb F_p|\mathbb F_0}$ et prendre
$a_i=0$ pour $1<i\le p-1$. Bien entendu, $F$ peut
être n'importe quel corps parfait \unsure{dans la sc.}
On peut aussi prendre $x,y$ alg. indép. sur $F$,
et $a_1$ alg. \deleted{\unsure{séparable}} \inserted{\unsure{irrationnel}} sur $F$, $a_1\notin F$, alors
$K$ est obtenu \unsure{en adjoignant} une racine
$p$-ième (celle de $y-ax$) dans $k(b)$, où $b$ est
alg. \unsure{séparable} sur $k$ $[b=a^{1/p}]$.

\mspage{28}
